\documentclass[11pt]{article}

\usepackage{amsmath,amssymb}
\usepackage{enumitem}
\usepackage{geometry}
\geometry{margin=1in}

\title{Mate's Immersed Boundary Method: Oral Exam Questions and Answers}
\author{}
\date{}

\begin{document}
\maketitle

\section*{IBM Mathematical Questions}

\begin{enumerate}[label=\arabic*.]

\item \textbf{What is the order of convergence of the immersed boundary method?}

\textit{Sample answer:}  
The immersed boundary method is first-order accurate globally. Away from the immersed boundary, the solution is second-order accurate pointwise. The reduction in global accuracy is due to the presence of singular forces at the immersed boundary.

\item \textbf{What is the order of convergence of the fluid solver you are using?}

\textit{Sample answer:}  
The underlying fluid solver is second-order accurate globally in space and time.

\item \textbf{Why does the overall IB method have lower accuracy than the fluid solver alone?}

\textit{Sample answer:}  
The immersed boundary introduces a singular force at the interface, which produces discontinuities in pressure and stress. These discontinuities are smoothed numerically using a regularized delta function, which degrades the global order of accuracy.

\item \textbf{Why not use fully implicit time stepping?}

\textit{Sample answer:}  
Fully implicit time stepping requires solving for the fluid velocity, pressure, and boundary configuration simultaneously. This leads to large, tightly coupled nonlinear systems. While implicit IB methods exist for linear forces (such as linear springs), extending them to nonlinear elastic forces is computationally expensive and often not beneficial in practice.

\item \textbf{How are you handling the singularity at the immersed boundary?}

\textit{Sample answer:}  
The singular force is regularized using a smooth delta function with finite width and height, which spreads the force over a small region of the grid.

\item \textbf{What are the key properties of the discretized delta function?}

\textit{Sample answer:}  
The discretized delta function is continuous, compactly supported, symmetric, conserves force and torque, and has self-coupling properties that are approximately independent of the relative position of Lagrangian points with respect to the Eulerian grid.

\item \textbf{What is the support of the discretized delta function?}

\textit{Sample answer:}  
For the standard 4-point delta function, the support extends to $\pm 2\Delta x$ in each coordinate direction. Other delta functions with smaller or larger support exist, but this choice is the most common.

\item \textbf{Is this the only discretized delta function you can use?}

\textit{Sample answer:}  
No. Other discretized delta functions exist, including kernels with wider support and simpler kernels such as the hat or cosine functions. The choice depends on accuracy, stability, and computational cost.

\item \textbf{Why not use a delta function with smaller support, such as one grid cell?}

\textit{Sample answer:}  
A support of two grid cells in each direction is the smallest support that can satisfy the desired moment and symmetry conditions. Kernels with smaller support can work in some cases, but may introduce artifacts or reduce stability.

\item \textbf{Why not use a delta function with higher-order moments?}

\textit{Sample answer:}  
Higher-order moment conditions require wider support, which increases computational cost and spreads forces over a larger region. The standard 4-point delta function represents a good compromise between accuracy and efficiency.

\item \textbf{If the delta function depends on grid spacing, why does the solution converge?}

\textit{Sample answer:}  
Although the regularized force depends on grid spacing, the Navier--Stokes equations integrate this force. As the grid is refined, the regularized force converges to a true singular force and the solution converges to the continuum IB formulation.

\item \textbf{Why does the interaction equation not require a Jacobian?}

\textit{Sample answer:}  
The force-spreading integral is not a change of variables but a distributional push-forward from material space to physical space. The integral is performed entirely over reference coordinates, so no Jacobian appears explicitly.

\end{enumerate}

\section*{IBM Limitations}

\begin{enumerate}[label=\arabic*.]

\item \textbf{Why are there no true boundary conditions in IBM?}

\textit{Sample answer:}  
Boundary conditions are enforced weakly. In the continuum limit, the interaction equations recover no-slip, a pressure jump proportional to the normal force, and a jump in the normal derivative of velocity proportional to the tangential force. This avoids the need for a body-fitted grid.

\item \textbf{Are you resolving pressure near the boundary?}

\textit{Sample answer:}  
No. The regularized delta smooths pressure near the boundary. However, the solution converges to the correct limit as the grid is refined.

\item \textbf{Are you resolving the boundary layer?}

\textit{Sample answer:}  
No. The boundary layer is smoothed by the delta function and is not sharply resolved.

\item \textbf{Can IBM resolve turbulent flows?}

\textit{Sample answer:}  
IBM is not well suited for resolving boundary layers in high-Reynolds-number turbulent flows. It can capture bulk flow features but not fine near-wall turbulence dynamics.

\item \textbf{Is IBM appropriate at low Reynolds number?}

\textit{Sample answer:}  
Yes. IBM can be used in the Stokes regime, although many researchers prefer the method of regularized Stokeslets, which exploits the linearity of the Stokes equations.

\item \textbf{You are not resolving the boundary layer. Justify the method.}

\textit{Sample answer:}  
IBM is appropriate when bulk flow dynamics are of interest and when resolving the immediate near-boundary structure is not critical. It is not suitable for applications that require detailed boundary-layer resolution.

\item \textbf{Does IBM have issues with volume conservation?}

\textit{Sample answer:}  
Yes. Classical IB methods can exhibit volume loss. More recent formulations address this by tuning discrete gradient and divergence operators to the interpolation scheme.

\end{enumerate}

\section*{Computational and Modeling Questions}

\begin{enumerate}[label=\arabic*.]

\item \textbf{Is your fluid solver implicit, explicit, or mixed?}

\textit{Sample answer:}  
It is a mixed method. The nonlinear convection and forcing terms are treated explicitly, while viscous terms are treated implicitly.

\item \textbf{What is the ratio of IB points to fluid grid points?}

\textit{Sample answer:}  
The spacing between IB points is approximately half the fluid grid spacing, corresponding to roughly two IB points per fluid grid cell.

\item \textbf{Why choose this ratio?}

\textit{Sample answer:}  
If IB points are too sparse, fluid leakage occurs. If they are too dense, numerical stiffness and spurious oscillations appear. This ratio has been found empirically to work well.

\item \textbf{Is the fluid discretized on a staggered grid?}

\textit{Sample answer:}  
No. In the FFT-based formulation used here, pressure--velocity decoupling is not an issue, so a staggered grid is not required.

\end{enumerate}

\section*{Other Methods}

\begin{enumerate}[label=\arabic*.]

\item \textbf{What are other approaches to fluid--structure interaction?}

\textit{Sample answer:}  
Other methods include the immersed interface method, ALE methods, level set methods, finite element and finite volume approaches, and the method of regularized Stokeslets.

\item \textbf{Why use IBM rather than these alternatives?}

\textit{Sample answer:}  
IBM avoids body-fitted meshes, handles large deformations naturally, and has robust open-source implementations. Many alternative methods are more difficult to implement or less suited for elastic structures.

\end{enumerate}

\end{document}
